% Input fragment; use the preamble of the English Volume II source.
% This is an ordinary chapter. Preserve existing special Chapter 5.5 counters.
\chapter{Proposed Chapter Title}\label{template:v2}
\begin{chaptersummary}
[AUTHOR TO COMPLETE: connect the established result, present question,
chosen structure or tool, later discretization use, and one core idea.]
\end{chaptersummary}
\studyroute{[Name omitted proofs, algebra, or implementation details and verified references.]}
{[Precise, established English search terms]}
{[Start from the current definitions or equations and request a checkable expansion.]}

\section{Objects, Conventions, and Question}
[Follow the volume's units, metric signature, and orientations.
Identify fields, spaces, and assumptions.]

% Optional draft author input: Volume II has no authornote environment.
% studentinput below is a reader exercise, not an author placeholder.
\paragraph{Potential Conceptual Jump}
\textbf{Potential question:} [Why introduce this object or operation here?]\par
\textbf{Author Supplement:} [AUTHOR TO COMPLETE: explain the motivation and selection rationale.]\par
\textbf{Core conceptual connection:} [Established result] $\longrightarrow$ [present question].

\section{Key Relations and Computation}
[Insert approved equations and explain boundary terms.
Identify continuous, semidiscrete, and fully discrete results.]

\begin{askai}
[Specify notation, assumptions, starting point, and a derivation or reproducible calculation.]

Follow-up: [Change one condition to test applicability.]
\end{askai}

% Optional active-recall and implementation blocks.
\begin{studentinput}
[State a specific task the reader can attempt independently.]

\textbf{Your answer:}\par\vspace{1cm}
\end{studentinput}
\begin{checkpoint}
[Ask a specific question testing a core conceptual distinction.]
\end{checkpoint}
\begin{engineernote}
[State the approved structure's implications for degrees of freedom,
interfaces, or implementation, with applicability conditions.]
\end{engineernote}

\section{Exercises}
\begin{enumerate}
\item {[Specify an exercise with explicit inputs and a checkable target.]}
\end{enumerate}
